\emph{Conventions.} Lattice weights, Haar measure, the absence of gauge fixing and $\beta = 2N/g^{2}$ ($= 6/g^{2}$ here) are those of Fact \claimref{synthesis.imported-wilson-lattice-gauge-theory-and-strong-coupling-confinement}(i); the tree gauge and the Haar change of variables are those of the proof of Theorem \claimref{synthesis.weak-factor-rotor-lattice-strong-coupling-area-law}(iv), (v). The colour factor is the $\mathfrak{su}(3)$ summand of Theorem \claimref{standard-model.gauge-algebra-minimal-massive-composite}, in the colour block of $\Cl(9,0)$ as in \claimref{synthesis.field-level-brst-differential-of-the-derived-gauge-fields-is-nilpotent}. A bivector $B$ acts on $V$ by $v \mapsto \tfrac12(vB - Bv)$; the check computes with $6 \times 6$ matrices of these actions and of signed permutations, each first matched with the Clifford product.

\emph{Step 1 (the colour group; computed, with a standard reading).} A rotation $U$ of $V$ commuting with $j$ is complex linear for the complex structure $j$ and unitary for the Hermitian form whose real part is the metric, so the group of such rotations is $U(3)$. On the complexified trivectors, $U$ preserves the complex line spanned by the $J$-volume form, and acts on it by its $J$-determinant (or its conjugate, by convention); on the real plane spanned by $\mathrm{Re}\,\Omega$ and $\mathrm{Im}\,\Omega$ this is the rotation by the argument of the determinant. So $U$ fixes $\mathrm{Re}\,\Omega$ exactly when its determinant is $1$, and the group is $SU(3)$, which is connected. The check confirms the convention on $\mathrm{diag}(j, 1, 1)$, of determinant $j$, which turns $\mathrm{Re}\,\Omega$ into $\mathrm{Im}\,\Omega$. Its Lie algebra consists of the bivectors $B$ of $V$ with $[B, J] = 0$ and $[B, \mathrm{Re}\,\Omega] = 0$ (the derivative of $RXR^{-1}$ along $R = \exp(\epsilon B/2)$ is $\tfrac12[B, X]$); solving these linear conditions on the fifteen bivectors of $V$ gives an eight-dimensional space, closed under the commutator and orthogonal to $J$ in the trace form: $e_{45} - e_{36}$, $e_{35} + e_{46}$, $e_{47} - e_{38}$, $e_{37} + e_{48}$, $e_{56} - e_{34}$, $e_{67} - e_{58}$, $e_{57} + e_{68}$, $e_{78} - e_{34}$. This is (i), first part.

\emph{Step 2 (the finite subgroup; computed).} Let $G$ be generated by the $j$-linear signed permutations with complex matrices $\mathrm{diag}(j, -j, 1)$, $\mathrm{diag}(1, j, -j)$, the cyclic permutation of $e_{3}, e_{5}, e_{7}$ (with $e_{4}, e_{6}, e_{8}$), and the exchange of the first two complex coordinates with the sign of the third reversed. The closure has $96$ elements, the monomial complex matrices with entries in $\{\pm1, \pm j\}$ and determinant $1$. For each element the check builds, by successive reflections, a product $R$ of an even number of vectors with rational components with $RvR^{-1} = g(v)$ on all six basis vectors, and confirms $RJR^{-1} = J$, $R\,\mathrm{Re}\,\Omega\,R^{-1} = \mathrm{Re}\,\Omega$ and $R\,\mathrm{Im}\,\Omega\,R^{-1} = \mathrm{Im}\,\Omega$; so $G \subset SU(3)$. An element of the form $a + bj$ acts on every complex coordinate by the same unit $a + bj$, and the only such signed permutations are $\pm1$, $\pm j$, with determinants $1$, $-1$, $\mp j$; so the centre of $G$ is trivial. The centre of $SU(3)$ consists of the $J$-scalars with $(a + bj)^{3} = 1$, that is $1$, $\omega$, $\omega^{2}$, $\omega = -\tfrac12 + \tfrac{\sqrt3}{2}j$; no element of $G$ has an irrational component, and the check confirms on a range of denominators that the only rational cube root of $1$ in $\{a + bj\}$ is $1$. This is (i).

\emph{Step 3 (averages are projectors).} Let $K$ be a compact group with normalised Haar measure acting orthogonally on a real space $W$ (here $W = V^{\otimes k}$, $U \mapsto U^{\otimes k}$), and $P = \int\rho(U)\,dU$. Invariance of the measure gives $\rho(h)P = P = P\rho(h)$, so $P^{2} = P$ and the image of $P$ is the invariant subspace; inversion invariance gives $P^{\mathsf T} = \int\rho(U^{-1})\,dU = P$. So $P$ is the orthogonal projector onto the $K$-invariants, and the same holds for the uniform average over a finite group. A polynomial of degree $k$ in the components of $U$ is a linear functional of $U^{\otimes k}$, so its average is determined by $P$. For $G \subset SU(3)$, the $SU(3)$-invariants are $G$-invariant; conversely, a $G$-invariant tensor annihilated by every element of $\mathfrak{su}(3)$ acting as a derivation is invariant under $\exp$ of the algebra, hence under the connected $SU(3)$. So if the $G$-invariants are all annihilated by $\mathfrak{su}(3)$, the two invariant subspaces coincide and the averages agree in that degree.

\emph{Step 4 (degrees one to three; proved with computation).} $G$ acts on the basis tensors of $V^{\otimes k}$ by signed permutations, so its invariants are spanned by the signed orbit sums of the orbits on which no element reverses a tensor's sign; the check finds $0$, $2$ and $2$ such orbits for $k = 1, 2, 3$, identifies the spans with the metric and the form of $J$ ($k = 2$) and with the tensors of $\mathrm{Re}\,\Omega$ and $\mathrm{Im}\,\Omega$ ($k = 3$), and finds that the eight derivations annihilate them. By Step 3 the averages agree in each degree $k \le 3$; the check also compares the $G$-average of $U^{\otimes k}$ entry by entry with the projector on a set of index pairs. For independent links the average of a product factorises (Fubini), and a polynomial of degree at most three in each link is a sum of products of per-link polynomials of degree at most three, so its averages agree as well. This is (ii).

\emph{Step 5 (degree four; computed).} For $k = 4$ the check finds $20$ signed orbits for $G$ and a $12$-dimensional subspace of their span annihilated by $\mathfrak{su}(3)$; so the invariant spaces differ, and $\int(\mathrm{tr}\,U)^{4}dU = \mathrm{tr}\,P = 12$ against $20$ for $G$ (also by summing over $G$). The derivation $D$ of $j$ commutes with $G$; it vanishes on the $12$ colour invariants, and on the $20$ invariants of $G$ its kernel has dimension $18$ and the kernel of $D^{2} + 16$ dimension $2$. On the complexification, $D$ is $j(p - q)$ on the tensors with $p$ factors from the $+j$ eigenspace and $q$ from the $-j$ eigenspace, and $\omega = \exp(2\pi J/3)$ acts there by $\omega^{p - q}$; so the two extra directions with $D^{2} = -16$ have $p - q = \pm4$, on which $\omega$ acts by $\omega^{\pm1} \neq 1$. Concretely, the $J$-complex trace $t = \sum_{k}(U_{2k,2k} + jU_{2k+1,2k})$ has $\tfrac{1}{96}\sum_{G}t^{4} = 1$ and $\tfrac{1}{96}\sum_{G}|t|^{4} = 3$ (computed), while for $SU(3)$ the substitution $U \mapsto \omega U$ preserves $dU$ and multiplies $t^{4}$ by $\omega^{4} = \omega$, so $\int t^{4}\,dU = 0$. This is (iii).

\emph{Step 6 (shared strong-coupling data; proved).} The gluing integrals are polynomials of degree one and two in $U$, so by Step 4 they may be computed over $G$; the check does so on the eighteen basis maps $E_{kl}$, $jE_{kl}$ of the $j$-linear maps (complete by bilinearity) and finds $0$ and $\tfrac13\mathrm{tr}(XY)$. Expanding the weight, the strong-coupling expansion of $\langle O\rangle$ for a gauge-invariant polynomial observable $O$ is a ratio of sums of integrals $\int O\prod_{p}\chi(U_{p})^{n_{p}}/n_{p}!$ over finitely many links, and each such integral with degree at most three in each link agrees by Step 4. For the rectangle, fix the tree of Theorem \claimref{synthesis.weak-factor-rotor-lattice-strong-coupling-area-law}(v) to $1$ (the Haar measure, and the average over $G$, are invariant under the change of variables): the plaquettes become $V_{x+1,y}V_{x,y}^{-1}$ and the loop the product of the right-hand column; integrating the columns one at a time with the gluing rule gives, at each of the $A = RT$ integrations, a factor $\tfrac13$ and a weight $\tfrac{\beta}{6}$, and the last trace $\tfrac16\mathrm{tr}\,1 = 1$, so the term is $(\tfrac{\beta}{6}\cdot\tfrac13)^{A} = (\beta/18)^{A}$. Every link of the rectangle occurs in two factors, degree two, so the same holds for $G$; the check sums over $G$ in the tree gauge for the $1 \times 1$ loop ($\tfrac{1}{18}$) and the $2 \times 1$ loop ($\tfrac{1}{324}$). A tube of plaquettes joining two plaquettes, each to first order, also has every link in exactly two factors. Whether $G$ has further terms below order $\beta^{A}$, from surfaces that use its absence of centre, is not decided here; for $SU(3)$ the centre excludes them. This is (iv).

\emph{Step 7 (the plaquette; proved).} With $m_{k}$ the moments of $\chi$, $m_{k} = 6^{-k}\int(\mathrm{tr}\,U)^{k}dU = 6^{-k}$ times the dimension of the invariants of $V^{\otimes k}$ (the trace of the projector, as $\mathrm{tr}(U^{\otimes k}) = (\mathrm{tr}\,U)^{k}$). For $SU(3)$: $m_{1} = 0$, $m_{2} = \tfrac{2}{36} = \tfrac1{18}$, $m_{3} = \tfrac{2}{216} = \tfrac1{108}$, $m_{4} = \tfrac{12}{1296} = \tfrac1{108}$; for $G$ the same except $m_{4} = \tfrac{20}{1296} = \tfrac{5}{324}$ (also by summing $\chi^{k}$ over $G$). The cumulants are $\kappa_{2} = m_{2}$, $\kappa_{3} = m_{3}$, $\kappa_{4} = m_{4} - 3m_{2}^{2}$, which is $0$ for $SU(3)$ and $\tfrac{1}{162}$ for $G$, and $(\ln z)' = \sum_{n}\kappa_{n+1}\beta^{n}/n!$ gives the two series; the check also expands the ratio $\int\chi e^{\beta\chi}/\int e^{\beta\chi}$ over $G$ directly. On the lattice, the character expansion $e^{\beta\chi} = z(\beta)\bigl(1 + \sum_{r \neq 1}d_{r}a_{r}(\beta)\chi_{r}\bigr)$ with $a_{r} = O(\beta)$ holds for any compact group; a product of non-trivial characters integrates to zero unless every link of its plaquettes lies in at least two of them (Schur, $\int\chi_{r}(XU)\,dU = 0$), and the smallest such set in dimension at least three is the boundary of an elementary cube, six plaquettes (standard, as in Drouffe and Zuber's account, cited from memory). So $\ln Z$ is the number of plaquettes times $\ln z$ plus terms of order $\beta^{6}$, and $\langle\chi(U_{p})\rangle = (\ln z)' + O(\beta^{5})$. This is (v).

\emph{Step 8 (the gap at the identity; computed).} The check takes the maximum of $\chi$ over the $95$ elements $g \neq 1$ of $G$: $\tfrac13$ (for instance at $\mathrm{diag}(j, -j, 1)$, real trace $2$). For $SU(3)$, $z = ((n^{2} - 1) + 2nj)/(n^{2} + 1)$ has $z\bar z = 1$, so $\mathrm{diag}(z, \bar z, 1) \in SU(3)$, with $\mathrm{tr} = 2 + 4(n^{2} - 1)/(n^{2} + 1)$ and $1 - \chi = 4/(3(n^{2} + 1))$, which tends to $0$. This is (vi).

\emph{Step 9 (reading; argued).} (a) By (ii) and (iv) no term of the strong-coupling expansion with link degree at most three separates $SU(3)$ from $G$, including the area-law terms $(\beta/18)^{A}$ and the tube terms of plaquette correlators from which the leading glueball decay rate is read. $G$ has no centre, so nothing like $N$-ality protects its fundamental loop at higher orders, and its weight is gapped at the identity by (vi), so at large $\beta$ each plaquette is frozen to $1$ except with weight at most $e^{-2\beta/3}$: a finite gauge group has no continuum limit, and by the standard low-temperature (Peierls) argument for finite gauge groups in four dimensions its Wilson loops have a perimeter law at weak coupling (cited from memory, not imported). The stored comparison with compact $U(1)$ gives the same form, $(\beta/2)^{A}$ (\claimref{synthesis.weak-factor-rotor-lattice-strong-coupling-area-law}(vi)), for a theory that deconfines with a massless phase at weak coupling (Fact \claimref{synthesis.imported-wilson-lattice-gauge-theory-and-strong-coupling-confinement}(iv)). Every finite subgroup of $SU(3)$ differs from Haar in some degree (a nonzero polynomial vanishes on any finite set), so such a double reaches only finitely many orders; but confinement and a mass gap in the continuum concern $\beta \to \infty$, reached by no finite order and needing asymptotic freedom ($b_{3} < 0$ under the stored content and premises, \claimref{synthesis.one-loop-running-coefficients-from-spacetime-algebra-spin-traces}). (b) Lattice spacing: the series' coefficients are completed to the reals (\claimref{synthesis.continuum-completion-premise}); P2's finiteness declaration is about the descent's deposited mode set (\claimref{dark-matter.p2-indifference}); the Planck cells of (C2) lie on the stretched singularity and are a Schwarzschild calibration (\claimref{singularity.c2-cell-density}); the common unit of (R2) is a tower Wigner-Seitz radius (\claimref{dark-matter.r2-one-quantum-per-unit-ball}); a finite floor of the tower keeps the rotations and no boost (\claimref{synthesis.floor-regulator-keeps-the-charge-and-the-central-term-below-the-cut-off}); the induced Newton coupling needs a cut-off the store does not fix (\claimref{synthesis.series-content-induces-a-newton-coupling-only-through-an-unfixed-cut-off}). None gives the gauge field a lattice spacing on flat spacetime. Were one adopted, the physical question would become one about a lattice at fixed spacing, where (iv) and Osterwalder and Seiler's results apply only at small $\beta$, and the mass-gap problem, which asks for the continuum theory on $\R^{4}$, would be untouched. (c) Gaps: the tower's pairing gap (\claimref{synthesis.tower-pairing-gap-at-fixed-charges-is-twice-the-quasiparticle-energy}) is a gap of fermionic quasiparticles of one floor at fixed charges; the doublet pair mode of a two-species floor (\claimref{synthesis.two-species-floor-pairing-has-a-gapped-doublet-pair-mode-that-closes-where-it-condenses}) is a composite of tower modes whose gap closes where it condenses; the window current algebra (\claimref{synthesis.exact-window-sugawara-identity-and-luttinger-normal-modes-of-tower-density-waves}) has density waves of energy $p$ at momentum $p$, gapless; free gluons are at null momenta (\claimref{synthesis.quartet-homotopy-makes-gluon-and-massive-vector-physical-states-positive-at-every-number-of-quanta}). None is a spectral gap of a Hamiltonian or transfer matrix of the four-dimensional gauge field.

\emph{Status.} (i), (iii), (vi) computed exactly; (ii), (iv), (v) proved, with the loops sampled and the closed-surface count of (v) standard; the reading argued. Standard lattice gauge theory recomputed in the colour block; transport-absent.
